\documentclass[../main.tex]{subfiles}

\begin{document}

\ifSubfilesClassLoaded{

    \tableofcontents
    
}{}
\section{Basic Topology}

\begin{problem}{Convex Bodies}{}
    A convex body is a compact, convex subset \(  K\subseteq \mathbb{R} ^{n}  \) which contains 0 as an interior point and which is symmetric around the origin .
    \begin{enumerate}
        \item Prove that the following formula defines a norm on  \(  \mathbb{R} ^{n}  \): \[
        N\left( x \right): =  \inf \left\{  \lambda > 0: \frac{x}{ \lambda }\in K \right\} 
        \] 
        \item Prove that \(  K  \) is homeomorphic to  \(  \mathbb{D}^{n}  \) and that \(   \partial K  \) is homeomorphicc to \(  \mathbb{S}^{n-1}  \)    . 
    \end{enumerate}
     
\end{problem}
\begin{note}
    \begin{enumerate}
        \item 若 \(  K  \)是一个单位球. 则 \(  N\left( x \right)= \left\| x \right\|   \).  
        \item 证明 \(  N\left( x \right)   \)是一个norm, 需要证明:1.非负正定 2.绝对一次齐次 3.三角不等式.    
        \item 由于 \(   \varphi   \)的形式足够简单, 可以通过直接构造逆映射证明双射.  
    \end{enumerate}
    
\end{note}

\begin{lemma}{}{}
    \(  X   \) is compact, \(  Y   \) is Hausdorff   . If \(  \exists f :X\to Y  \) is a continious bijection, then \(  X,Y  \) are homeomorphic.   
\end{lemma}
\begin{proof}
    \begin{enumerate}
        \item It is obvious that \(  N\left( x \right)\ge 0   \).  If \(  N\left( x \right)= 0   \) , then \[
        \inf \left\{  \lambda > 0: \frac{x }{ \lambda  }\in K  \right\}= 0
        \]   which means that for all  \(   \varepsilon > 0  \),   \(  \frac{x }{  \varepsilon   } \in K   \) .Since \(  K  \) is compact, \(  \mathrm{dist}\left( K,0 \right)< \infty   \).      If \(  x \neq 0  \) ,then \(  \left\| x \right\|> 0  \) . We take   \(   \varepsilon _0   < \frac{1 }{\mathrm{dist}\left( K,0 \right) \cdot \left\| x \right\| } \) . Then \(  \left\| \frac{x }{ \varepsilon _0  }  \right\|>   \mathrm{dist}\left( K,0 \right) \)    . But \(  \frac{x }{  \varepsilon _0  }    \in K\), which is a contradiction.For the second part, let \(  x \in \mathbb{R} ^{n}, k \in \mathbb{R} \setminus \left\{ 0 \right\}  \). For all \(   \lambda > N\left( x \right)   \) ,  \[
        \frac{x }{ \lambda  }\in K \implies  \frac{kx }{k \lambda  } \in K  
        \]  ,which shows that \(  N\left( kx \right)\le kN\left( x \right)    \)    .  Replace \(  k   \) by \(  \frac{1 }{ k}   \) and replace \(  x  \) by \(  kx  \), we have \(  N\left( kx \right)\ge kN\left( x \right)    \). Thus \(  N\left( kx \right)= kN\left( x \right)    \).       For the triangle inequality , we need to show that \[
       \inf \left\{  \lambda > 0:\frac{x_1+ x_2 }{ \lambda  }\in K  \right\} \le \inf \left\{  \lambda > 0: \frac{x_1 }{ \lambda  }\in K  \right\}+ \inf \left\{  \lambda > 0: \frac{x_2 }{ \lambda  }  \in K\right\}
        \] Let \(   \lambda _0 = N\left( x_1+ x_2 \right)   \) .  Only need to show that \[
        \frac{x_1+ x_2 }{N\left( x_1 \right)+ N\left( x_2 \right)   } \in K 
        \]For all \(   \delta > 0  \), we have  \[
        \frac{x_1 }{N\left( x_1 \right)+  \delta   }\in K,\quad \frac{x_2 }{N\left( x_2 \right)+  \delta   }  \in K
        \]. Since \(  K  \) is convex , we have the convex combination \[
        \frac{1 }{N\left( x_1 \right)+ N\left( x_2 \right)+ 2 \delta    } \left( \left( N\left( x_2 \right)+  \delta   \right)\frac{x_2 }{N\left( x_2 \right)+  \delta   }+ \left( N\left( x_1 \right)+  \delta   \right)\frac{x_1 }{N\left( x_1 \right)+  \delta   }     \right) \in K 
        \]That is \[
        \frac{x_1+ x_2 }{N\left( x_1 \right)+ N\left( x_2 \right)+ 2 \delta    }\in K 
        \] for all \(   \delta > 0  \). Then \[
        N\left( x_1+ x_2 \right)\le N\left( x_1 \right)+ N\left( x_2 \right)+ 2 \delta ,\quad \forall  \delta > 0 \implies N\left( x_1+ x_2 \right)\le N\left( x_1 \right)+ N\left( x_2 \right)      
        \] 
        \item    \(  x\in K^{\circ }\iff  N\left( x \right)< 1   \) ,  \(  x \in \mathrm{Ext}\left( K \right)\iff  N\left( x \right)> 1    \) ,  \(  x \in  \partial K  \iff N\left( x \right)= 1 \).    We define a map \(   \varphi : \mathbb{R} ^{n}\to \mathbb{R} ^{n}  \) . \[
         \varphi \left( x \right): =\begin{cases}  N\left( x \right) \frac{x }{\left\| x \right\| }  , &  x\neq 0\\ 0,& x = 0   \end{cases} 
        \]We have  \[
        \left\|  \varphi \left( x \right)  \right\|= N\left( x \right)
        \]   Then \(   \varphi   \) maps  \(  K^{\circ}  \) on to  \(  \left( \mathbb{D} ^{n}\right)^{\circ}    \), and maps \(   \partial K  \) on to \(  \mathbb{S}^{n-1}  \).  To show \(   \varphi   \) is continious , only need to show it is continious at \(  0  \) .  In fact,  \(\left\|  \lim_{x\to 0} \varphi \left( x \right)  \right\|=   \lim_{x\to 0}\left\|  \varphi \left( x \right)  \right\|=\lim_{x\to 0} N\left( x \right)= 0   \) since \(  N  \) is a norm   . We define  \(  \psi : \mathbb{R} ^{n}\to \mathbb{R} ^{n}  \) : \[
        \psi \left( x \right):= \begin{cases} \frac{\left\| x \right\| }{N\left( x \right)  }x ,& x \neq 0\\ 
         0,& x= 0  \end{cases}  
        \]   Then \(   \varphi \circ \psi = \psi \circ  \varphi = \operatorname{Id}  \) .   The same skills can be used to show \(  \psi   \) is continious at \(  0  \)  as well . 
        Finally , by the lemma above, \(  X,Y  \) are homeomorphic.  
    \end{enumerate}
    
\end{proof}


\begin{problem}{Projective spaces}{}
    Let \(  \mathbb{K}= \mathbb{R}  \) or \(  \mathbb{C}  \) . The \(   n  \) -dimensional projectinve space \(  \mathbb{KP}^{n}  \) is the orbit space of \(  \mathbb{k}^{n+ 1}\setminus \left\{ 0 \right\}   \) under the action of \(  \mathbb{k}^{*}  \) by rescaling. 
    \begin{enumerate}
        \item Prove that \(  \mathbb{RP}^{n}  \) is homeomorphic to the orbit space \(  \mathbb{S}^{n} \setminus \left\{ \pm 1 \right\}  \) , where \(  \left\{ \pm 1 \right\}   \) acts by multiplication on \(  \mathbb{S}^{n}\subseteq \mathbb{R} ^{n+ 1}  \) . 
        \item Prove that \(  \mathbb{CP}^{n}  \) is homeomorphic to the orbit space \(  \mathbb{S}^{2n+ 1} \setminus  \mathbb{S}^{1}  \) , where \(  \mathbb{S}^{1}= \left\{ z \in \mathbb{C}  : \left| z \right|= 1 \right\}  \) acts on \(  \mathbb{S}^{2n+ 1} \subseteq \mathbb{C}^{n+ 1}  \) by multiplication. 
        \item Prove that \(  \mathbb{RP}^{1}  \) is homeomorphic to \(  \mathbb{S}^{1}  \) and that \(  \mathbb{CP}^{1}  \) is homeomorphic to \(  \mathbb{S}^{2}  \) .             
    \end{enumerate}
           
\end{problem}
\begin{proof}
    \begin{enumerate}
        \item Define \[
         \pi: \mathbb{R} ^{n+ 1}\setminus \left\{ 0 \right\}\to \mathbb{RP}^{n} ,\quad  \pi\left( x \right)= \left[ x \right]  
        \] and define \[
        f: \mathbb{S}^{n}\to \mathbb{RP}^{n}=   \left.  \pi \right|_{\mathbb{S}^{n}}
        \] Then \(  f  \) is continuous.  Take any \(  p \in \mathbb{RP}^{n}  \), there exists \(  v \in \mathbb{R} ^{n+ 1}\setminus \left\{ 0 \right\}  \), such that \(  p= \left[ v \right]   \).  Let \(  x = \frac{v }{\left\| v \right\| }   \), then \(  x \in \mathbb{S}^{n}  \)     , \[
        f\left( x \right)= \left[ x \right]= \left[ v \right]= p   
        \] thus \(  f  \) is surjective.  If \[
        f\left( x \right)= f\left( y \right)  
        \] , then \(  \left[ x \right]= \left[ y \right]    \), there exists \(   \lambda  \in \mathbb{R}^{*}  \), such taht \(  x=  \lambda y  \). Since \[
        \left\| x \right\|= \left\|  \lambda y \right\|= \left|  \lambda  \right|\left\| y \right\|\implies \left|  \lambda  \right|= 1  
        \], \(   \lambda = \pm 1  \), \(  x= y  \) or \(  x = -y  \).Thus \(  \mathrm{Fiber}_{p}\left( f \right)= \mathrm{Orbit}_{x}\left( \left\{ \pm 1 \right\} \right)    \)  The universal property of quotient topology  gives a continuous bijection \[
        \tilde{f}: \mathbb{S}^{n}\setminus \left\{ \pm 1 \right\}\to \mathbb{RP}^{n}
        \]    Since \(  \mathbb{S}^{n}\setminus \left\{ \pm 1 \right\}  \) is compact, \(  \mathbb{RP}^{n}  \) is Hausdorff,  \(  \tilde{f}  \) is a homeomorphism.    
        \item Similarly, let \(  g:\mathbb{S}^{2n+ 1}\to \mathbb{CP}^{n}  \), \(  g\left( z \right)= \left[ z \right]    \). Take any \(  p = \left[ w \right] \in \mathbb{CP}^{n}  \) ,where \(  w \in \mathbb{C}^{n+ 1}\setminus \left\{ 0 \right\}  \). Let \(  z = \frac{w }{\left\| w \right\| }   \). Then \[
        g\left( z \right)= \left[ z \right]= \left[ w \right]= p   
        \] \(  g  \) is surjective.      
        If \[
        g\left( z_1 \right)= g\left( z_2 \right)  
        \] , then there exists \(   \lambda  \in \mathbb{C}^{*}  \), such that \(  z_1=  \lambda z_2  \), then \[
        \left\| z_1 \right\|= \left|  \lambda  \right|\left\| z_2 \right\|\implies \left|  \lambda  \right|= 1  
        \] which gives that \(   \lambda \in \mathbb{S}^{1}  \). Thus \(  \mathrm{Fiber}_{p}\left( f \right)\simeq \mathrm{Orbit}_{x}\left( \mathbb{S}^{1} \right)    \).  Thus \[
        \mathbb{CP}^{n}\simeq  \mathbb{S}^{2n+ 1}/ \mathbb{S}^{1}
        \]    
        \item \begin{enumerate}
            \item Only need to shwo that \(  \mathbb{S}^{1}/\left\{ \pm 1 \right\}\simeq \mathbb{S}^{1}  \), \[
            h\left( z \right)= z^{2} 
            \] induces a continuous bijection \(  \tilde{h}: \mathbb{S}^{1}\setminus \left\{\pm 1 \right\}\to \mathbb{S}^{1}  \). \(  \mathbb{S}^{1}\setminus \left\{ \pm 1 \right\}  \) is compact, \(  \mathbb{S}^{1}  \) is Hausdorff, \(  \tilde{h}  \) is a homemophism.
            \item Since \(  \mathbb{S}^{3}/\mathbb{S}^{1}\simeq \mathbb{S}^{2}  \).        
        \end{enumerate}
        
    \end{enumerate}
    
\end{proof}
\begin{problem}{}{}
    
The torus $\mathbb{T}$ is the quotient of $\left[ 0,1 \right] ^2$ by the equivalence relation generated by $(x,0) \sim (x,1)$ and $(0,y) \sim (1,y)$ for all $x, y \in \left[ 0,1 \right]^{2} $. Prove that $\mathbb{T}$ is homeomorphic to:
\begin{enumerate}
    \item The product $\mathbb{S}^1 \times \mathbb{S}^1$.
    \item The quotient of $\mathbb{R}^2$ under the action of the (discrete) group $\mathbb{Z}^2$ by translations.
\end{enumerate}
\end{problem}
\begin{proof}
    \begin{enumerate}
        \item Let \[
        f\left( x,y \right):=  \left( e^{2 \pi xi}, e^{2 \pi yi} \right)  
        \] Then \[
        f\left( x,0 \right)= \left( e^{2 \pi i x},1 \right)= f\left( x,1 \right)   
        \] \[
        f\left( 0,y \right)= \left( 1, e^{2 \pi yi} \right)= f\left( 1,y \right)   
        \]Thus \(  f  \) induces a map \(  \bar{f}  \) on \(  \mathbb{T}  \) such that \(  \bar{f}\left( [\left( x,y \right) ] \right)=\left( e^{2 \pi xi} , e^{2 \pi yi}\right)     \) , which is a continuous function. It is not hard to see that \(  \bar{f}: \mathbb{T}\to \mathbb{S}^{1}\times \mathbb{S}^{1}  \) is a bijection.  As a continuous image of \(  \left[ 0,1 \right]^{2}   \), \(  \mathbb{T}  \) is compact.   Take any closed subset \(  V  \) of \(  \mathbb{T}  \), \(  V  \) is compact. \(  \bar{f}\left( V \right)    \) is compact as the continuous   image of \(  V  \). Since \(  \mathbb{S}^{1}\times \mathbb{S}^{1}  \) is Hausdorff, \(  \bar{f}\left( V \right)   \) is closed.   Thus \(  \bar{f}  \) is a continuous and closed bijection, thus is a homeomorphism. 
        \item Let \[
        f\left( x,y \right)= \left( x,y \right)/ \mathbb{Z}^{2} 
        \] Then Since \(  \left( x,0 \right)-\left( x,1 \right)\in \mathbb{Z} ^{2}, \left( 0,y \right)-\left( 1,y \right)\in \mathbb{Z} ^{2}      \). \(  f  \) induces a continuous map \(  \bar{f}  \)  on from \(  \mathbb{T}  \) to \(  \mathbb{S}^{1}\times \mathbb{S}^{1}  \)   .   
        It is clear that \(  f  \) maps onto \(  \mathbb{R} ^{2}/\mathbb{Z}^{2}  \). If \(  f\left( x_1,y_1 \right)= f\left( x_2,y_2 \right)    \) , then \[
        \left( x_1-x_2,y_1-y_2 \right)\in \mathbb{Z} ^{2} \implies  x_1= x_2 \text{ or } \left\{ x_1,x_2 \right\}= \left\{ 0,1 \right\}, \text{and }y_1= y_2\text{ or } \left\{ y_1,y_2 \right\}= \left\{ 0,1 \right\}
        \]  thus \(  \left( x_1,y_1 \right)\sim \left( x_2,y_2 \right)    \) .  Again since \(  \mathbb{T}  \) is compact, \(  \mathbb{R} ^{2}/ \mathbb{Z} ^{2}  \) is Hausdorff, \(  \bar{f}  \) is a homemorphism.      
    \end{enumerate}
    
\end{proof}

\begin{problem}{}{}
    Consider the orthogonal group $O_{n-1}(\mathbb{R})$ as the subgroup of $O_n(\mathbb{R})$ of matrices of the form:
\[
\begin{pmatrix}
1 & 0 \\
0 & A
\end{pmatrix} \in O_n(\mathbb{R}), A \in O_{n-1}(\mathbb{R}).
\]
Let $O_{n-1}(\mathbb{R})$ act on $O_n(\mathbb{R})$ by left multiplication. Prove that the orbit space $O_n(\mathbb{R})/O_{n-1}(\mathbb{R})$ is homeomorphic to $\mathbb{S}^{n-1}$.
\end{problem}
\begin{proofsketch}
    将 \(  O_{n}\left( \mathbb{R}  \right)   \)刻画为 \(  \mathbb{S}^{n-1}  \)上的群作用. 此时 \(  A \sim \begin{pmatrix} 
        1&0\\ 
         0&A 
    \end{pmatrix}   \)就是北极点的稳定子群. 一但 作用是传递的, 由流形上的轨道-稳定化子定理, 就有   \[
    O_{n}\left( \mathbb{R}  \right)/ O_{n-1}\left( \mathbb{R}  \right)\simeq \mathbb{S}^{n-1}  
    \]
\end{proofsketch}


\section{Homotopy}

\begin{problem}{Mobius band}{}
    Let $M$ be the M\"obius band, that is, the quotient of $[0, 1]^2$ under the equivalence relation generated by $(x,0) \sim (1 - x, 1)$ for all $x \in [0,1]$. Prove that $M$ is homotopy equivalent to $\mathbb{S}^1$. 
\end{problem}
\begin{proof}
    Define a map \[
    H: \left[ 0,1 \right]^{2}\times \left[ 0,1 \right] \to \mathbb{S}^{1},  \quad H\left( \left( x,y \right),t  \right)= \left( \left( 1-t \right)x+ \frac{1}{2}t,y  \right) 
    \]   Let \(  \tilde{H}=  \pi\circ H  \). Then \[
    \tilde{H}\left( \left( x,0 \right),t  \right)=  \pi\left(  \left( 1-t \right)x+ \frac{1}{2}t,0  \right)  
    \] \[
    \tilde{H}\left( \left( 1-x ,1\right),t  \right)=   \pi\left( \left( \left( 1-t \right)\left( 1-x \right)+ \frac{1}{2}t   \right),1  \right)  
    \]  Since \[
  1-  \left( \left( 1-t \right)x+ \frac{1}{2}t  \right)= \left( 1-t \right)\left( 1-x \right)+ \frac{1}{2}t   
    \], \[
    \tilde{H}\left( \left( x,0 \right),t  \right)= \tilde{H}\left( \left( 1-x,1 \right),t  \right)  
    \] \(  \tilde{H}  \) induces a continuous map \(  H^{*}  \)  on \(  M  \).   And \[
  H^{*}\left( \left[ \left( x,y \right),0  \right]  \right)= \left[ \left( x,y \right)  \right],\quad H^{*}\left( \left[ \left( x,y \right),1  \right]  \right)= \left[ \left( \frac{1}{2},y \right)  \right]    
    \]Thus \(  H^{*}  \) is a deformation restract frome \(  M  \) to \(  \mathbb{S}^{1}  \), which is a homotopy equivalence as well.    
\end{proof}
\begin{problem}{}{}
    Let X be space. The cone CX is the quotient space:
\[
CX = (X \times [0,1])/(X \times \{0\}).
\]
Let $\pi: X \times [0,1] \to CX$ be the quotient map and $\iota : X \to CX, x \mapsto \pi(x,1)$.
\begin{enumerate}
    \item Let $f: X \to Y$ be a map. Prove that $f$ is homotopic to a constant map if and only if there exists $f': CX \to Y$ such that $f' \circ \iota = f$.
    \item Prove that $CS^n$ is homeomorphic to $D^{n+1}$.
    \item Prove that for all X, CX is contractible.
\end{enumerate}
\end{problem}
\begin{proof}
    \begin{enumerate}
        \item If such \(  f^{\prime}   \) exists. Then \(  H: =  f^{\prime} \circ  \pi  : X\times \left[ 0,1 \right]\to Y \) is a homotopy from \(  f  \) to the constant map valued the vetex of \(  CX  \). Conversly, if \(  f  \) is homotopic to constant map. Suppose \(  H:X\times I\to Y  \) is its homotopy map.  Since \(  CX  \) is the quotient space of \(  X\times \left[ 0,1 \right]   \), and \(  H\left( X\times \left\{ 0 \right\} \right)=  \text{constant}   \). It follows frome the universal property that  \(  H  \) induces a map \(  f^{\prime} :CX\to Y  \) such that \(  f^{\prime} \circ  \pi= H  \). Finally \[
       \left(  f^{\prime} \circ \iota  \right)\left( x \right)= f^{\prime} \left(  \pi\left( x,1 \right)  \right)= H\left( x,1 \right)= f\left( x \right)     
        \]        
        \item Define \[
        g: \mathbb{S}^{n}\times \left[ 0,1 \right]\to D^{n+ 1},\quad g\left( x,t \right)= t\cdot x  
        \] Since \(  g\left( x,0 \right)\equiv 0\in D^{n+ 1}   \).  \(  g  \) induces a  continuous function \(  h: CS^{n}\to D^{n+ 1}  \) such that \(  g\circ \pi= h  \). We now show that \(  h  \) is a bijection.  If  \[
        h\left( \left[ x_1,t_1 \right]  \right)= h\left( \left[ x_2,t_2 \right]  \right)  
        \], that is \(  t_1x_1= t_2x_2  \) , then \[
        \left\| t_1x_1 \right\|= \left\| t_2x_2 \right\|\implies  t_1= t_2
        \] If \(  t_1= t_2  = 0\) , then \(  \left[ x_1,t_1 \right]= \left[ x_2,t_2 \right]= 0    \).  If \(  t_1= t_2> 0  \) , then  \(  x_1= x_2  \), thus \(  \left[ x_1,t_1 \right]= \left[ \times 2,t_2 \right]    \) . Thus \(  h  \) is injective. For surjectivity, let \(  y \in D^{n+ 1}  \). If \(  y= 0  \) , we take \(  t= 0  \), then \(  h\left( \left[ x,0 \right]= 0  \right)   = y\). If \(  y \neq 0  \) , then \[
        h\left( \left[ \frac{y }{\left\| y \right\| },\left\| y \right\|  \right]  \right)= y 
        \]  Thus \(  h  \) is surjective. Since \(  CS^{n}  \) is compact, \(  D^{n+ 1}   \)            is Hausdorff, \(  h  \) is homemorphism.  
        \item Define \[
        H:CX \times \left[0,1 \right]\to CX,\quad H\left( \left[ x,s \right],t  \right)= \left[ x,s\cdot \left( 1-t \right)  \right]   
        \] Then it is not hard to see that \(  H  \) is well-defined and continuous. Since \[
        H\left( \left[ x,s \right],0  \right)= \left[ x,s \right],\quad H\left( \left[ x,s \right],1  \right)= \left[ x,0 \right]    
        \] \(  H  \) is a homotopy from \(  \operatorname{Id}   \) to \(  \text{Constant}  \)   .
    \end{enumerate}
    
\end{proof}
\begin{problem}{Linear Groups}{}
   
\begin{enumerate}
    \item Prove that the inclusion $O_n(\mathbb{R}) \to GL_n(\mathbb{R})$ is a homotopy equivalence. (Use the Gram-Schmidt orthonormalization procedure to construct the reverse map).
    \item Among the following matrix groups, determine which ones are compact, and determine their $\pi_0$:
    $GL_n(\mathbb{C})$, $GL_n(\mathbb{R})$, $O_n(\mathbb{R})$, $SO_n(\mathbb{R})$, $U_n(\mathbb{C})$, $SU_n(\mathbb{C})$.
\end{enumerate}
\end{problem}
\begin{proof}
    \begin{enumerate}
        \item Take \(  A \in \operatorname{GL} _{n}\left( \mathbb{R}  \right)   \), suppose \(  A= \left(  v_1,\cdots,v_n  \right)   \). Then \(  \left\{  v_1,\cdots,v_n  \right\}  \) passing to an orthonormal basis \(  \left\{  u_1,\cdots,u_n  \right\}  \) by Gram-Schmidt orthonormalization. We write \(  r\left( A \right)= \left\{  u_1,\cdots,u_n  \right\}   \). Then \(  r  \) is a continuous map on \(  \operatorname{GL} _{n}\left( \mathbb{R}  \right)   \) since the inner product is continuous. We define \[
        H:\operatorname{GL} _{n}\left( \mathbb{R}  \right)\times \left[ 0,1 \right]\to \operatorname{GL} _{n}\left( \mathbb{R}  \right),\quad H\left( A,t \right)= \left( 1-t \right)A+ tr\left( A \right)      
        \] we need to verify that \(  H\left( A,t \right) \in \operatorname{GL} _{n}\left( \mathbb{R}  \right)    \) for all \(  A,t  \)         . The \(  k  \)-th column vector of \(  H\left( A,t \right)   \) is \(  \left( 1-t \right)v_{k}+ tu_{k}   \). Since \[
       v_{k} \in \operatorname{span}\left\{ u_1,\cdots ,u_{k-1} \right\}
        \] we can write \[
        \left( 1-t \right)v_{k}+ tu_{k}=  \left( 1-t \right)\sum _{j= 1}^{k}c_{kj}u_{j}+ tu_{k}  
        \]   Since \(  c_{jj}> 0  \), \(  \left( 1-t \right)c_{j j}+ t > 0  \), \(  \left\{ \left( 1-t \right)v_1+ tu_1,\cdots ,\left( 1-t \right)v_{k}+ tu_{k}   \right\}  \) are linealy independent.   And since \[
        H\left( A,0 \right)= A= \operatorname{Id}_{\operatorname{GL} _{n}\left( \mathbb{R}  \right) }\left( A \right)  
        \] \[
        H\left( A,1 \right)= r\left( A \right) \in O_{n}\left( \mathbb{R}  \right)   
        \] and for \(  B \in O_{n}\left( R \right)   \) ,\(  H\left( B,t \right)= r\left( B \right)= 1    \).  We have \(  H:\operatorname{Id}_{\operatorname{GL} _{n}\left( \mathbb{R}  \right) }\simeq i  \) is a SDR.
        \item \(  \operatorname{GL} _{n}\left( \mathbb{C}  \right)   \) is a open submanifold of \(  M_{n}\left( \mathbb{C}  \right)   \), which is noncompact. Since \(  \operatorname{GL} _{n}\left( \mathbb{C}  \right)   \) is path-connected. \(   \pi_0\left( \operatorname{GL} _{n}\left( \mathbb{C}  \right)  \right)=   1 \)  . \(  \operatorname{GL} _{n}\left( \mathbb{R}  \right)   \) is a open submanifold of \(  M_{n}\left( \mathbb{R}  \right)   \), thus \(  Gl_{n}\left( \mathbb{R}  \right)   \) is noncompact. \(  \det : \operatorname{GL} _{n}\left( \mathbb{R}  \right)\to \mathbb{R} \setminus \left\{ 0 \right\}   \) is a continuous surjection.  Since \(  \mathbb{R} \setminus \left\{ 0 \right\}   \) has two components. \(  \operatorname{GL} _{n}\left( \mathbb{R}  \right)   \) has at least two components.  Since \[
        \operatorname{GL} _{n}^{+ }\left( \mathbb{R}  \right)= \left\{ A: \det A> 0 \right\},\quad \operatorname{GL} _{n}^{-}\left( \mathbb{R}  \right)= \left\{ A: \det A< 0 \right\}  
        \] are exactly two components of \(  \operatorname{GL} _{n}\left( \mathbb{R}  \right)   \). \(   \pi_0\left( \operatorname{GL} _{n}\left( \mathbb{R}  \right)  \right)= 2   \). 
         Since \(  \left\{ A: A^{T}A= E \right\}  \) is closed.  \(  O_{n}\left( \mathbb{R}  \right)   \) is bounded.  \(  O_{n}\left( \mathbb{R}  \right)   \) is compact. \(   \pi_0\left( O_{n}\left( \mathbb{R}  \right)  \right)= 2   \).   \(  SO_{n}\left( \mathbb{R}  \right)   \) is a closed subset of \(  O_{n}\left( \mathbb{R}  \right)   \), thus is compact as well. \(   \pi_0\left( SO_{n}\left( \mathbb{R}  \right)  \right)= 1   \)              . 
         \(  U_{n}\left( \mathbb{C}  \right)   \) is a closed subset of \(  \operatorname{GL} _{n}\left( \mathbb{C}  \right)   \). \(   \pi_0\left( U_{n}\left( \mathbb{C}  \right)  \right)= 1   \)    . \(  SU_{n}\left( \mathbb{C}  \right)   \) is compact, \(   \pi_0\left( SU_{n}\left( \mathbb{C}  \right)  \right)= 1   \).   
    \end{enumerate}
    
\end{proof}
\begin{problem}{}{}
    \begin{enumerate}
    \item Prove that $\mathbb{CP}^{n+1}$ is obtained from $\mathbb{CP}^n$ by gluing a cell of dimension $2n+2$.
    \item Compute the fundamental group of $\mathbb{CP}^n$ for $n \geq 1$.
\end{enumerate}
\end{problem}
\section{van Kampen}

\begin{lemma}{}{}
    设 \(  X  \)分解为一些道路连通集 \(  A_{\alpha }  \)的并,且每个 \(  A_{\alpha }  \)包含了基点 \(  x_0 \in X  \).则所有含入映射 \(  A_{\alpha }\hookrightarrow X  \)的诱导同态 \(  j_{\alpha }:\pi _1 \left( A_{\alpha } \right)\to \pi _1 \left( X \right)    \)扩张为同态 \(  \Phi : *_{\alpha } \pi _1 \left( A_{\alpha } \right)\to \pi _1 \left( X \right)    \).      
\end{lemma}

\begin{proposition}{}{}
    令 \(  i_{\alpha \beta }:\pi _1 \left( A_{\alpha }\cap A_{\beta } \right)\to \pi _1 \left( A_{\alpha } \right)    \)表示含入映射 \(  A_{\alpha }\cap A_{\beta }\hookrightarrow A_{\alpha }  \)诱导的同态,则 \(  j_{\alpha }i_{\alpha \beta }= j_{\beta }i_{\beta \alpha }  \).所以 \(  \operatorname{ker}\,\Phi   \)包含所有形如 \(  i_{\alpha \beta }\left(  \omega  \right)i_{\beta \alpha }\left(  \omega  \right)^{-1} , \omega \in \pi _1 \left( A_{\alpha }\cap A_{\beta } \right)     \)     的元素.
\end{proposition}
\begin{note}
    即 \(  \Phi   \)将 \(  A_{\alpha }\cap A_{\beta }  \)落在\(  A_{\alpha }  \)和 \(  A_{\beta }  \)    中部分的像等同.
\end{note}

\begin{theorem}{van Kampen}{}
    设 \(  X  \)分解为包含了基点 \(  x_0 \in X  \)的道路连通开集 \(  A_{\alpha }  \)的并.且每个交集 \(  A_{\alpha }\cap A_{\beta }  \)是道路连通的,则 同态 \(  \Phi : *_{\alpha }\pi _1 \left( A_{\alpha } \right)\to \pi _1 \left( X \right)    \)  是满射. 此外,若每个 \(  A_{\alpha }\cap A_{\beta }\cap A_{ \gamma }  \) 也是道路连通的,则 \(  \Phi   \)的核 \(  N  \)由全体形如 \(  i_{\alpha \beta }\left(  \omega  \right)i_{\beta \alpha }\left(  \omega  \right)^{-1} , \omega \in \pi _1 \left( A_{\alpha }\cap A_{\beta } \right)     \)的元素生成,从而 \(  \Phi   \)诱导出同构\[
    \pi _1 \left( X \right)\simeq *_{\alpha }\pi _1 \left( A_{\alpha } \right) / N=  *_{\alpha }\pi _1 \left( A_{\alpha } \right)/ \left( i_{\alpha \beta }\left(  \omega  \right)= i_{\beta \alpha }\left(  \omega  \right)   \right)    
    \]    
\end{theorem}

\begin{example}{}{}
    设 \(  \left\{ X_{\alpha } \right\}  \)是一族道路连通的空间, \(  x_{\alpha } \in X_{\alpha }  \)是基点,考虑它们的楔和 \(\bigvee  _{\alpha }X_{\alpha }  \).    对于每个 \(  x_{\alpha }  \)  ,它在 \(  X_{\alpha }  \)中都存在形变收缩到自身的开邻域 \(  U_{\alpha }  \).令 \(  A_{\alpha }= X_{\alpha } \bigvee _{\beta \neq \alpha }U_{\beta }  \),则 \(  A_{\alpha }  \)形变收缩到 \(  X_{\alpha }  \)  .此外, \(  \bigcap_{\alpha }A_{\alpha }   \)形变收缩到点 \(  x_{\alpha }  \).  由van Kampen定理,  \(  \Phi :\pi _1 \left( X \right)\to *_{\alpha }\pi _1 \left( A_{\alpha } \right)\simeq *_{\alpha }\pi _1 \left( X_{\alpha } \right)     \)  是同构.
\end{example}

\hspace*{\fill} 

\section{Homology}




\begin{lemma}{Snake lemma}{}
Given a commutative diagram of $R$-module homomorphisms: where the two horizontal sequences are exact, there exists a $R$-module homomorphism

\[\begin{tikzcd}
	& {M'} & M & {M''} & 0 \\
	0 & {N'} & N & {N''}
	\arrow["\alpha", from=1-2, to=1-3]
	\arrow["{f'}", from=1-2, to=2-2]
	\arrow["\beta", from=1-3, to=1-4]
	\arrow["f", from=1-3, to=2-3]
	\arrow[from=1-4, to=1-5]
	\arrow["{f''}", from=1-4, to=2-4]
	\arrow[from=2-1, to=2-2]
	\arrow["{\alpha'}", from=2-2, to=2-3]
	\arrow["{\beta'}", from=2-3, to=2-4]
\end{tikzcd}\]


$\delta : \text{Ker } f'' \longrightarrow \text{Coker } f'$, called the \emph{connecting homomorphism} such that the sequence
$$
\text{Ker } f' \xrightarrow{\bar{\alpha}} \text{Ker } f \xrightarrow{\bar{\beta}} \text{Ker } f'' \xrightarrow{\delta} \text{Coker } f' \xrightarrow{\bar{\alpha'}} \text{Coker } f\xrightarrow{\bar{\beta'}} \text{Coker } f''
$$
is exact. Moreover, the \dfntxt{connecting homomorphism} $\delta$ has the naturality properties, so that the above assignment of a ‘snake’ to the corresponding ‘six-term’ exact sequence of modules defines a covariant functor.
\end{lemma}

\begin{note}
	最关键的信息是 \dfntxt{connetcting homomorphism} \(   \delta   \).  构造的过程就是借助 \(  \beta   \)的是surjective, \(  \alpha ^{\prime}   \) 是 injectinve, 来调整使得 \(  \beta   \)和 \(  \alpha ^{\prime}   \)对应的箭头是可以反转的. 事实上在应用时, 这个connetcting homomorphism的具体构造往往比它的存在性更重要.    
\end{note}

\begin{proof}\dfntxt{to be finished}
	We need to show
	\begin{enumerate}
		\item \(  \bar{\alpha},\bar{\beta}  \) are well defined, where \(  \bar{\alpha}:=  \alpha |_{\operatorname{ker}f^{\prime} }, \bar{\beta}: = \beta |_{\operatorname{ker}f}  \) 
		\item \(  \operatorname{Im}\bar{\alpha}= \operatorname{ker}\bar{\beta}  \) .
		\item \(  \bar{\alpha}^{\prime} , \bar{\beta}^{\prime}   \) are well defined, where \(  \bar{\alpha}^{\prime} \left( [x] \right):= [\alpha ^{\prime} \left( x \right) ]   \) , \(  \bar{\beta}^{\prime} \left( [y] \right):= [\beta ^{\prime} \left( y \right) ]   \)   
		\item \(  \operatorname{Im} \bar{\alpha}^{\prime} =\operatorname{ker}\bar{\beta}^{\prime}   \) 
		\item Construct a  \(   \delta   \).  
		\item \(  \operatorname{Im}\bar{\beta}  = \operatorname{ker} \delta \)
		\item \(  \operatorname{Im} \delta = \operatorname{ker}\bar{\alpha}^{\prime}   \)  
	\end{enumerate}
	We  prove the above points below.
	\begin{enumerate}
		\item   By commutative, \(  f\circ \alpha \left( \operatorname{ker}f^{\prime}  \right)=  \alpha ^{\prime} \circ f^{\prime} \left( \operatorname{ker}f^{\prime}  \right)= 0    \), which shows that  \(  \alpha \left( \operatorname{ker}f^{\prime}  \right)\subseteq \operatorname{ker}f   \). Thus \(  \bar{\alpha}  \) is well defined. \(  \bar{\beta}  \)'s is similar.
		\item    \[
		\operatorname{Im}\bar{\alpha}= \alpha \left( \operatorname{ker}f^{\prime}  \right)\subseteq \operatorname{Im}\alpha = \operatorname{ker}\beta \implies  \beta \left( \operatorname{Im}\bar{\alpha} \right)= 0 \implies \operatorname{Im}\bar{\alpha}\subseteq \operatorname{ker}\bar{\beta}  
		\] The other side, \[
		\operatorname{ker}\bar{\beta}= \operatorname{ker}\beta \cap \operatorname{ker}f= \operatorname{Im}\alpha \cap \operatorname{ker}f
		\] Take \(  y \in \operatorname{ker}\bar{\beta}  \), there exists \(  x \in M^{\prime}   \) such that \(  y= \alpha \left( x \right)   \) and  \(  f\left( y \right)= f\left( \alpha \left( x \right)  \right)= 0    \). Since \(  f\circ \alpha = \alpha ^{\prime} \circ f^{\prime}   \), \(  \alpha ^{\prime} \circ f^{\prime} \left( x \right)= 0   \). Since \(  \alpha ^{\prime}   \) is injective, \(  x \in \operatorname{ker}f^{\prime}   \). Thus \(  y \in \operatorname{Im}\bar{\alpha}  \).          
		\item For \(  \bar{\alpha}^{\prime}   \) , we need to show that for \(  x_1,x_2 \in N^{\prime}  \) such that \(  x_1-x_2 \in \operatorname{Im}f^{\prime}   \), \(  \alpha ^{\prime} \left( x_1 \right)-\alpha ^{\prime} \left( x_2 \right)\in \operatorname{Im}f    \)   . It is true by commutative \[
		  \alpha ^{\prime} \left( x_1 \right)-\alpha ^{\prime} \left( x_2 \right)= \alpha ^{\prime} \left( x_1-x_2 \right) \in \operatorname{Im}\left( \alpha ^{\prime} \circ f^{\prime}  \right)= \operatorname{Im}\left( f\circ \alpha  \right)   \subseteq \operatorname{Im}f    
		\] \(  \bar{\beta}^{\prime}   \)'s is similar.
		\item  Take \(  [y]_{\sim \operatorname{Im}f } \in \operatorname{Im}\bar{\alpha}^{\prime}   \), then  there exists \(  [x]_{\sim \operatorname{Im}f^{\prime}  } \in \operatorname{Coker}f^{\prime}   \), such that \(  \alpha ^{\prime} \left( [x] _{\sim \operatorname{Im}f^{\prime} }\right)= [y]  _{\sim \operatorname{Im}f} \). Then \(  \alpha ^{\prime} \left( x \right)-y \in \operatorname{Im}f   \), there exists \(  z \in M  \), such that  \[
		\alpha ^{\prime} \left( x \right)-y= f\left( z \right)  
		\]  By commutative and exactness: \(  \operatorname{ker}\beta ^{\prime} = \operatorname{Im}\alpha ^{\prime}   \),   \[
		\beta ^{\prime} \left( y \right)= \beta ^{\prime} \left( \alpha ^{\prime} \left( x \right)-f\left( z \right)   \right)=0- \beta ^{\prime} \circ f\left( z \right)       =- f^{\prime \prime} \circ \beta \left( z \right) \in \operatorname{Im}f^{\prime \prime}  
		\], which shows that  \(  \bar{\beta}^{\prime} \left( [y]_{\sim \operatorname{Im}f} \right)   = [0]_{\sim \operatorname{Im}f^{\prime \prime} }\), \(  \operatorname{Im}\bar{\alpha}^{\prime}   \subseteq \operatorname{ker}\bar{\beta}^{\prime} \). 
		
		Then, take \(  [y]_{\sim \operatorname{Im}f}\in \operatorname{ker} \bar{\beta}^{\prime}   \). Then \(  y\in \operatorname{ker}\beta ^{\prime} = \operatorname{Im}\alpha ^{\prime}   \). There exists \(  x \in N^{\prime}   \) such that   \(  y= \alpha ^{\prime} \left( x \right)   \).  Then \(  [y]_{\sim \operatorname{Im}f}= \alpha ^{\bar{\prime}}\left( [x]_{\sim \operatorname{Im}f^{\prime} } \right)   \), \(  [y]_{\sim \operatorname{Im}f}\in \operatorname{Im} \bar{\alpha}^{\prime}   \), \(  \operatorname{ker}\bar{\beta}^{\prime} \subseteq \operatorname{Im} \bar{\alpha}^{\prime}  \).   
		\item Take \(  z \in \operatorname{ker}f^{\prime \prime}   \). Since \(  \beta   \) is surjective, there exists \(  y \in M  \) such that \(  \beta \left( y \right)= z   \). Once we show that  \(   f\left( y  \right)\in \operatorname{Im}\alpha ^{\prime}    \)       , we can define\[
	 \delta \left( z \right)= [\left({\alpha}^{\prime}  \right)^{-1} \left( f\left( y\right)  \right)  ]_{\sim \operatorname{Im}f^{\prime} }
	\] provided that this definition is well defined. By commutative \[
	\beta ^{\prime} \left( f\left( y \right)  \right)= f^{\prime \prime} \left( \beta \left( y \right)  \right)= f^{\prime \prime} \left( z \right)= 0 .  
	\] Hence \(   f\left( y \right) \in \operatorname{ker}\beta ^{\prime} = \operatorname{Im} \alpha ^{\prime}   \). Finally, to show that it is well defined, we take \(  z_1,z_2\in \operatorname{ker}f^{\prime \prime}   \) such that  \(  \beta \left( y_1 \right)= \beta \left( y_2 \right)= z    \). We need to show that \( \left( \alpha ^{\prime}  \right)^{-1} \left( f\left( y_1-y_2 \right)    \right)\in \operatorname{Im}f^{\prime}      \).  Since \(  \beta \left( y_1 \right)-\beta \left( y_2 \right)= 0    \), \(  \beta \left( y_1-y_2 \right)= 0   \), \(  y_1-y_2\in \operatorname{ker}\beta = \operatorname{Im}\alpha   \). Suppose \(  \alpha \left( x \right)= y_1-y_2   \). Then \[
	\left( \alpha ^{\prime}  \right)^{-1} \left( f\left( y_1-y_2 \right)  \right)= \left( \alpha ^{\prime}  \right)^{-1} \left( f\alpha \left( x \right)  \right)= \left( \alpha ^{\prime}  \right)^{-1} \left( \alpha ^{\prime} \circ f^{\prime}  \right)\left( x \right)= f^{\prime} \left( x \right) \in \operatorname{Im}f^{\prime}         
	\]    , which completes the proof.
	\end{enumerate}
	
	
\end{proof}

\begin{corollary}{}{}
Consider the following commutative diagram of $R$-modules and $R$-linear maps in which the two rows are exact. If $f_1$ and $f_3$ are isomorphisms then so is $f_2$.

\centering
\begin{tikzcd}
0 \arrow[r] & M_1 \arrow[r, "\alpha_1"] \arrow[d, "f_1"] & M_2 \arrow[r, "\alpha_2"] \arrow[d, "f_2"] & M_3 \arrow[r] \arrow[d, "f_3"] & 0 \\
0 \arrow[r] & N_1 \arrow[r, "\beta_1"] & N_2 \arrow[r, "\beta_2"] & N_3 \arrow[r] & 0
\end{tikzcd}
\end{corollary}
\begin{proof}
	By snake lemma, we have \[
	0\xrightarrow{\alpha _1 } \operatorname{ker}f_2\xrightarrow{\alpha _2 }0\xrightarrow{ \delta }0\xrightarrow{\beta _1 }\operatorname{Coker}f_2\xrightarrow{\beta _2 }0
	\]is exact. Then we have \[
	\begin{aligned}
	0&=  \operatorname{ker}\alpha _2 = \operatorname{Im}\alpha _1 = 0\\ 
	 &=  \operatorname{ker}f_2
	\end{aligned}
	\] which shows that \(  \operatorname{ker}f_2= 0  \). Similarly, \(  \operatorname{Coker} f_2  = 0\). The above shows that \(  f_2  \) is an isomorphisim.  
\end{proof}
\begin{corollary}{Four lemma}{}
Consider the following commutative diagram of $R$ modules and $R$-linear maps in which the two rows are exact. Suppose that $f_1$ is surjective and $f_4$ is injective. Then
\begin{enumerate}
    \item[(i)] $f_2$ is injective $\implies f_3$ is injective.
    \item[(ii)] $f_3$ is surjective $\implies f_2$ is surjective.
\end{enumerate}
\begin{center}
\begin{tikzcd}
M_1 \arrow[r, "\alpha_1"] \arrow[d, "f_1"] & M_2 \arrow[r, "\alpha_2"] \arrow[d, "f_2"] & M_3 \arrow[r, "\alpha_3"] \arrow[d, "f_3"] & M_4 \arrow[d, "f_4"] \\
N_1 \arrow[r, "\beta_1"] & N_2 \arrow[r, "\beta_2"] & N_3 \arrow[r, "\beta_3"] & N_4
\end{tikzcd}
\end{center}

\end{corollary}

\begin{proofsketch}
	如果 \(  f_3  \)被两个单射 \(  f_2,f_4  \)夹在中间,得益于 \(  f_1  \)是满射, 可以将右三列保持单射性地调整为一条snake, 从而导出 \(  f_3  \)是单射.
	
	类似地,如果 \(  f_2  \)被两个满射 \(  f_1,f_3  \)夹在中间, 得益于 \(  f_4  \)是单射, 左三列可以保持满射性地调整为一条snake, 导出 \(  f_2  \)是满射.  
\end{proofsketch}

\begin{proof}
\begin{enumerate}
	\item 
We can construct a snake \[\begin{tikzcd}
	& {M_2/\operatorname{ker}\alpha _2 } & {M_3} & {\operatorname{Im}\alpha _3} & 0 \\
	0 & {N_2/ \ker{\beta_2}} & {N_3} & {N_4}
	\arrow["{{\alpha_2}}", from=1-2, to=1-3]
	\arrow["{{\bar{f_2}}}", from=1-2, to=2-2]
	\arrow["{{\alpha_3}}", from=1-3, to=1-4]
	\arrow["{{f_3}}", from=1-3, to=2-3]
	\arrow[from=1-4, to=1-5]
	\arrow["{{ f_4|_{\operatorname{Im}\alpha _3 }  }}", from=1-4, to=2-4]
	\arrow[from=2-1, to=2-2]
	\arrow["{{\beta_2}}", from=2-2, to=2-3]
	\arrow["{{\beta_3}}", from=2-3, to=2-4]
\end{tikzcd}\] By snake lemma, the following sequence is exact \[
  \operatorname{ker}\bar{f}_{2}\to \operatorname{ker}f_3\to \operatorname{ker}f_4|_{\operatorname{Im}\alpha _3 }
\] Since \(  f_4  \) is injective,  \(  f_4|_{\operatorname{Im}_{\alpha _3 }}  \) is as well. Furthermore, by applying the snake lemma to the following diagram \[\begin{tikzcd}
	& {M_1} & {M_2} & {M_2/\ker \alpha_2} & 0 \\
	0 & {N_1/\ker{\beta_1}} & {N_2} & {N_2/\ker{\beta_2}}
	\arrow["{{{\alpha_1}}}", from=1-2, to=1-3]
	\arrow["{{{\bar{f_1}}}}"', from=1-2, to=2-2]
	\arrow["p", from=1-3, to=1-4]
	\arrow["{{{f_2}}}", from=1-3, to=2-3]
	\arrow[from=1-4, to=1-5]
	\arrow["{{{\bar{f_2}}}}", from=1-4, to=2-4]
	\arrow[from=2-1, to=2-2]
	\arrow["{{{\bar{\beta_1}}}}"', from=2-2, to=2-3]
	\arrow["p"', from=2-3, to=2-4]
\end{tikzcd}\] We get \[
0= \operatorname{ker}f_2\to \operatorname{ker}\bar{f_2}\xrightarrow{ \delta } \operatorname{Coker}\bar{f_1}
\]is a exact suquence,where \(  \operatorname{Coker}\bar{f_1}= 0  \) since \(  f_1  \) is surjective and \(  \bar{f_1}  \) is as well.   It shows that \(  \operatorname{ker} \bar{f}_{2}= 0  \) . Finally, we have the following sequence is exact \[
0= \operatorname{ker}\bar{f}_{2} \to \operatorname{ker}f_3 \to \operatorname{ker}f_4|_{\operatorname{Im}\alpha _3 }= 0
\]It follows that \(  \operatorname{ker}f_3= 0  \), \(  f_3  \) is injective.   
\item Another snake we can construct is \[\begin{tikzcd}
	& {M_1} & {M_2} & {\operatorname{Im}\alpha _2= \operatorname{ker}\alpha _3  } & 0 \\
	0 & {N_1/\ker{\beta_1}} & {N_2} & {\operatorname{Im}\beta _2 = \operatorname{ker}\beta_3  }
	\arrow["{\alpha_1}", from=1-2, to=1-3]
	\arrow["{\bar{f_1}}", from=1-2, to=2-2]
	\arrow["{\alpha_2}", from=1-3, to=1-4]
	\arrow["{f_2}", from=1-3, to=2-3]
	\arrow[from=1-4, to=1-5]
	\arrow["{f_3|_{\operatorname{Im}\alpha _2 } }", from=1-4, to=2-4]
	\arrow[from=2-1, to=2-2]
	\arrow["{\bar{\beta_1}}", from=2-2, to=2-3]
	\arrow["{\beta_2}", from=2-3, to=2-4]
\end{tikzcd}\]Then by snake lemma ,the following sequence is exact \[
\operatorname{ker}\bar{f_1}\to \operatorname{ker}f_2\to \operatorname{ker}f_3|_{\operatorname{Im}\alpha _2 }\to \operatorname{Coker}\bar{f_1}\to \operatorname{Coker}f_2\to \operatorname{Coker}f_3|_{\operatorname{Im}\alpha _2 }
\] Furthermore, by applying the snake lemma to the following diagram 
\[
\begin{tikzcd}[column sep=large, row sep=large]
	0 \arrow[r] & \operatorname{Ker}\alpha_3 \arrow[r, "\iota_M"] \arrow[d, "f_3|_{\operatorname{Im}\alpha_2}"] & M_3 \arrow[r, "\alpha_3"] \arrow[d, "f_3"] & \operatorname{Im}\alpha_3 \arrow[r] \arrow[d, "{f_4|_{\operatorname{Im}\alpha_3}}"] & 0 \\
	0 \arrow[r] & \operatorname{Ker}\beta_3 \arrow[r, "\iota_N"] & N_3 \arrow[r, "\beta_3"] & \operatorname{Im}\beta_3 \arrow[r] & 0
\end{tikzcd}
\] We get \[
0= \operatorname{ker}f_4|_{\operatorname{Im}\alpha _3 }\xrightarrow{ \delta }\operatorname{Coker}\left( f_3|_{\operatorname{Im}\alpha _2 } \right)\to \operatorname{Coker}f_3= 0 
\]which shows that \(  \operatorname{Coker}\left( f_3|_{\operatorname{Im}\alpha _2 } \right) = 0  \). Finally, we have the following sequence exact \[
0= \operatorname{Coker}\bar{f_1}\to \operatorname{Coker}f_2\to \operatorname{Coker}f_3|_{\operatorname{Im}\alpha _2 }= 0
\] Thus \(  \operatorname{Coker}f_2= 0  \), \(  f_2   \) is surjective.  
\end{enumerate}
   
\end{proof}


\begin{corollary}{Five lemma}{}
In the following diagram of $R$-modules, the two rows are given to be exact. If $f_1, f_2, f_4$ and $f_5$ are isomorphisms then $f_3$ is also an isomorphism.

\centering
\begin{tikzcd}
M_1 \arrow[r] \arrow[d, "f_1"] & M_2 \arrow[r] \arrow[d, "f_2"] & M_3 \arrow[r] \arrow[d, "f_3"] & M_4 \arrow[r] \arrow[d, "f_4"] & M_5 \arrow[d, "f_5"] \\
M_1' \arrow[r] & M_2' \arrow[r] & M_3' \arrow[r] & M_4' \arrow[r] & M_5'
\end{tikzcd}
\end{corollary}


\begin{proof}
	By applying the Four lemma to the left four columns, we get \(  f_3  \) is injective. And by applying the Four lemma to the right four columns , we get \(  f_3  \) is surjective .  

\end{proof}



\begin{theorem}{}{7-9-2}
Given a short exact sequence of chain complexes
\[
0 \to C'_* \xrightarrow{\alpha} C_* \xrightarrow{\beta} C''_* \to 0
\]
there is a functorial long exact sequence of homology groups
\[
\to H_n(C'_*) \xrightarrow{H_n(\alpha)} H_n(C_*) \xrightarrow{H_n(\beta)} H_n(C''_*) \xrightarrow{\delta_n} H_{n-1}(C'_*) \xrightarrow{H_{n-1}(\alpha)} H_{n-1}(C_*) \to
\]
\end{theorem}
\begin{proofsketch}
	Consider the diagram 

% https://q.uiver.app/#q=WzAsOCxbMSwwLCJDJ19uL1xcbWF0aHJte0ltfSBcXHBhcnRpYWwnX3tuKzF9Il0sWzIsMCwiQ19uL1xcbWF0aHJte0ltfSBcXHBhcnRpYWxfe24rMX0iXSxbMywwLCJDJydfbi9cXG1hdGhybXtJbX0gXFxwYXJ0aWFsJydfe24rMX0iXSxbNCwwLCIwIl0sWzAsMSwiMCJdLFsxLDEsIlxcbWF0aHJte0tlcn0gXFxwYXJ0aWFsJ197bi0xfSJdLFsyLDEsIlxcbWF0aHJte0tlcn0gXFxwYXJ0aWFsX3tuLTF9Il0sWzMsMSwiXFxtYXRocm17S2VyfSBcXHBhcnRpYWwnJ197bi0xfSJdLFswLDEsIlxcYmFye1xcYWxwaGF9X24iXSxbMCw1LCJcXHBhcnRpYWwnX24iXSxbMSwyLCJcXGJhcntcXGJldGF9X24iXSxbMSw2LCJcXHBhcnRpYWxfbiJdLFsyLDNdLFsyLDcsIlxccGFydGlhbCcnX3tuKzF9Il0sWzQsNV0sWzUsNiwiXFxhbHBoYSdfe24tMX0iXSxbNiw3LCJcXGJldGEnX3tuLTF9Il1d
\[\begin{tikzcd}
	& {C'_n/\mathrm{Im} \partial'_{n+1}} & {C_n/\mathrm{Im} \partial_{n+1}} & {C''_n/\mathrm{Im} \partial''_{n+1}} & 0 \\
	0 & {\mathrm{Ker} \partial'_{n-1}} & {\mathrm{Ker} \partial_{n-1}} & {\mathrm{Ker} \partial''_{n-1}}
	\arrow["{\bar{\alpha}_n}", from=1-2, to=1-3]
	\arrow["{\partial'_n}", from=1-2, to=2-2]
	\arrow["{\bar{\beta}_n}", from=1-3, to=1-4]
	\arrow["{\partial_n}", from=1-3, to=2-3]
	\arrow[from=1-4, to=1-5]
	\arrow["{\partial''_{n+1}}", from=1-4, to=2-4]
	\arrow[from=2-1, to=2-2]
	\arrow["{\alpha'_{n-1}}", from=2-2, to=2-3]
	\arrow["{\beta'_{n-1}}", from=2-3, to=2-4]
\end{tikzcd}\]
\end{proofsketch}
\begin{example}{reduced homology}{}
    定义 \(   \varepsilon   \)为将0-链映为系数和的同态. 称同态 \(  \varepsilon   \)为增广同态,可以按以下方式扩张出链复形 \(  \tilde{S}\left( X \right)    \)           \[
         \tilde{S}_n(X)=\left\{\begin{array}{ll}S_n(X),&\mathrm{~if~}n\geq0,\\\mathbb{Z},&\mathrm{~if~}n=-1,\\(0),&\mathrm{~if~}n<-1;\end{array}\right.\quad\mathrm{~and~}\quad\tilde{\partial}_n=\left\{\begin{array}{ll}\partial_n,&n\geq1,\\\varepsilon,&n=0,\\0,&n<0.\end{array}\right.
     \]相应的同调群记作 \(  \tilde{H}_{*}\left( X \right)   \),称为约化同调群.显然 \(  \tilde{H}_{i}\left( X \right)= \left( 0 \right),i<0;\tilde{H_{i}}\left( X \right)\simeq H_{i}\left( X \right), i>1; \tilde{H_0}\left( X \right) \oplus \mathbb{Z} \simeq H_0\left( X \right)        \)  .特别地,对于道路连通空间 \(  X  \), \(  \tilde{H_0}\left( X \right)= 0   \).  
 \end{example}
 
 \begin{problem}{cofibration}{}
    
Let $(X, A)$ be a pair of spaces. We say that the inclusion $A \to X$ is a cofibration if, whenever $f : X \to Y$ and $H : A \times [0,1] \to Y$ are maps such that $H(a, 0) = f (a)$ for all $a \in A$, there exists $\tilde{H} : X \times [0,1] \to Y$ such that $\tilde{H}(a, t) = H(a, t)$ for all $(a, t) \in A \times [0,1]$ and $\tilde{H}(x, 0) = f (x)$ for all $x \in X$.
\begin{enumerate}
    \item Prove that the inclusion $(X, A) \to (X \cup CA, CA)$ induces an isomorphism on relative homology.
    \item Prove that if $A \to X$ is a cofibration and $A$ is contractible, then the quotient map $X \to X/A$ is a homotopy equivalence.
    \item Prove that if $X$ is obtained from $A$ by gluing cells, then $A \to X$ is a cofibration.
    \item Prove that if $A \to X$ is a cofibration, then $(X \cup CA, CA)$ is a cofibration.
    \item Let $A \to X$ be a cofibration. Prove that the quotient map induces an isomorphism $H_i(X,A) \cong H_i(X/A,A/A) = \tilde{H}_i(X/A)$.
\end{enumerate}
 \end{problem}

 

\begin{problem}{homology of suspension}{}
    Let $X$ be a space. Compute $H_*(\Sigma X)$ in terms of $H_*(X)$.

\end{problem}
\begin{problem}{}{}
    Compute the homology of $\mathbb{CP}^n$.
\end{problem}
\section{Degree}
\begin{problem}{}{}
    
Let $f : S^n \to S^n$ be a map. Its \dfntxt{degree} $\deg(f) \in \mathbb{Z}$ is the integer such that for all $x \in H_n(S^n; \mathbb{Z})$, we have $f_*(x) = \deg(f) \cdot x$.
\begin{enumerate}
    \item Prove that this definition matches with the definition with fundamental groups for $n = 1$.
    \item Prove that if $\deg(f) \neq 0$ then $f$ is surjective. Find a counterexample for the converse.
    \item Prove that if $f$ is injective, then $\deg(f) = \pm 1$. Find a counterexample for the converse.
    \item Prove that the degree of a reflection is $-1$ (use Mayer-Vietoris). Given $A \in O_n(\mathbb{R})$, prove that the degree of $x \mapsto Ax$ is equal to $\det(A)$. Compute the degree of the antipodal map $x \mapsto -x$.
    \item Prove that $\Sigma S^n$ is homeomorphic to $S^{n+1}$, then prove that the degree of $f$ is equal to the degree of its suspension $\Sigma f : S^{n+1} \to S^{n+1}$.
\end{enumerate}
\end{problem}
\begin{problem}{}{}
    
\begin{enumerate}
    \item Prove that any map $f : S^n \to S^n$ without fixpoints is homotopic to the antipodal map. (Use that for all $x \in S^n$, the origin $0 \in \mathbb{R}^{n+1}$ is not on the segment $[f(x), -x]$).
    \item Prove that if $n$ is even, then any $f : S^n \to S^n$ admits a fixpoint.
    \item Prove that if a group $G$ acts freely on $S^n$ with $n$ even, then $\#G \le 2$.
    \item Prove that $S^n$, for $n$ odd, admits a nonvanishing vector field.
    \item Prove that any vector field on $S^n$, for $n$ even, vanishes at some point. (Given such a vector field $X$, consider the map $u \mapsto X(u)/\|X(u)\|)$.
\end{enumerate}

\end{problem}
\end{document}
